\section{Source curation}
\subsection{Detailed implementation}
\label{appendix:source_curation}

For tokens within a span, we compute weighted votes that combine:
(i) part-of-speech–dependent semantic weights,
(ii) confidence scores $c_t$ (optionally exponentiated),
and (iii) run-length–based coherence signals that favor temporally contiguous source assignments.

Formally, the curation function
$$
\hat{\mathcal{S}}_k
=
\texttt{CurateSourcesWithConf}
$$
$$
\big(
\{\hat{s}(t)\}_{t\in T_k}, \\
\{\texttt{pos}_t\}_{t\in T_k},
\{c_t\}_{t\in T_k},
\texttt{cfg}
\big)
$$
selects the minimal set of source units whose normalized attribution mass satisfies coverage and stability constraints. This procedure filters spurious token-level fluctuations, enforces cross-token consistency, and yields concise span-level explanations. .

Together, these three stages transform arbitrary token-level attribution signals into stable, generation-aware span-level source explanations, satisfying the design requirements outlined in \Cref{sec:problem}.

\begin{lstlisting}
def curate_sources_with_conf(
    source_ids: List[int],
    pos: List[str],
    conf: List[float],
    cfg: SourceCurationConfig = DEFAULT_CURATE_CFG,
) -> List[int]:
    T = len(source_ids)
    if T == 0:
        return []
    if not (T == len(pos) == len(conf)):
        raise ValueError("source_ids, pos, conf must have the same length.")

    vote: List[float] = []
    for p, c in zip(pos, conf):
        pw = POS_W.get(p, 0.3)
        cw = (max(c, 0.0) ** cfg.gamma)
        vote.append(pw * cw)

    total = float(sum(vote))
    if total <= 0:
        return []

    mass = defaultdict(float)
    for s, v in zip(source_ids, vote):
        mass[s] += v
    p_mass = {s: m / total for s, m in mass.items()}

    run_max = defaultdict(float)
    cur_s, cur_run = source_ids[0], vote[0]
    for i in range(1, T):
        if source_ids[i] == cur_s:
            cur_run += vote[i]
        else:
            run_max[cur_s] = max(run_max[cur_s], cur_run)
            cur_s, cur_run = source_ids[i], vote[i]
    run_max[cur_s] = max(run_max[cur_s], cur_run)
    run_frac = {s: run_max[s] / total for s in mass.keys()}

    def score(s: int) -> float:
        return cfg.alpha * p_mass[s] + (1.0 - cfg.alpha) * run_frac[s]

    ranked = sorted(p_mass.keys(), key=score, reverse=True)

    selected: List[int] = []
    cum = 0.0
    for s in ranked:
        strong_run = run_frac[s] >= cfg.run_min
        if p_mass[s] < cfg.p_min and not strong_run:
            continue
        selected.append(s)
        cum += p_mass[s]
        if cum >= cfg.coverage:
            break

    return selected

\end{lstlisting}

\subsection{Ablation study results on Audio and Video Tasks}
\label{appendix:ablation_filtering_audio_video}

Here we report the ablation results on Audio and Video tasks in addition to visual tasks in \Cref{tab:ablation_filtering}.

\begin{table}[h]
\centering
\tablesize
\caption{
Ablation study of the source curation framework. 
We report Precision, Recall, and F1 for audio and video attribution on the 134 audio and 102 video QA tasks.
The full model integrates POS-aware weighting and run-level coherence..
}
\label{tab:ablation_filtering_audio_video}
\skiptablebegin
\setlength{\tabcolsep}{4pt}
\begin{tabular}{l|ccc|ccc}
\toprule
\multirow{2}{*}{\textbf{Method}} 
& \multicolumn{3}{c|}{\textbf{Audio Attribution}} 
& \multicolumn{3}{c}{\textbf{Video Attribution}} \\
\cmidrule(lr){2-4} \cmidrule(lr){5-7}
& Precision & Recall & F1 
& Precision & Recall & F1 \\
\midrule
Full Model (Default) & 65.96 & 47.27 & 49.90 & 55.57 & 47.08 & 40.16 \\
\midrule
w/o POS Weighting & 74.66 & 43.94 & 50.07 & 56.85 & 41.30 & 37.46 \\
w/o Confidence Weight & 70.95 & 46.26 & 50.83 & 56.36 & 43.50 & 38.82 \\
w/o Confidence      & 74.26 & 42.33 & 48.69 & 56.48 & 38.21 & 35.80 \\
w/o Run Coherence & 74.26 & 42.23 & 48.69 & 56.48 & 38.20 & 35.79 \\
w/o $p_{\min}$ Filtering & 74.21 & 42.49 & 48.85 & 56.62 & 38.77 & 36.22 \\
\bottomrule
\end{tabular}
\skiptableend
\end{table}

\subsection{Ablation study}
\label{appendix:ablation_filtering}
Below, we provide a detailed analysis of each ablation setting, clarifying the functional role of each component in the source curation pipeline and its empirical impact.

\paragraph{Role of POS-aware weighting.}
POS-aware weighting assigns higher contribution to semantically informative tokens (e.g., nouns, proper nouns, numerals) and downweights function words and punctuation during attribution mass aggregation.
This mechanism ensures that source selection is driven primarily by content-bearing tokens rather than syntactic scaffolding.
Disabling POS-aware weighting produces only minor changes in text, audio and video attribution. However, image F1 decreases drastically to 20.79.
This indicates that visual grounding is particularly sensitive to semantic token emphasis: without syntactic weighting, low-information tokens contribute equally,
causing attribution mass to disperse across irrelevant image sources and destabilizing cross-modal assignments.

\paragraph{Role of confidence shaping.}
Confidence shaping modulates token contributions according to the strength of their token-to-source linkage, amplifying high-certainty alignments while attenuating weak or ambiguous ones.
Concretely, the exponent parameter $\gamma$ sharpens attribution mass toward confident mappings.
Without confidence modulation, attribution mass becomes diffuse across competing candidates, leading to over-selection of spurious image sources.
The larger degradation in image F1 suggests that visual grounding signals are inherently sparser and noisier than textual alignments,
making confidence-aware filtering critical in cross-modal scenarios.

\paragraph{Role of run-level coherence.}
Run-level coherence promotes temporal consistency by favoring source chunks that receive sustained support across consecutive generated tokens.
This mechanism captures the intuition that semantically grounded evidence often corresponds to contiguous spans (e.g., multiple tokens describing the same object).
Without enforcing temporal consistency, scattered token-level assignments accumulate across unrelated image sources, undermining stable visual grounding.
These results indicate that coherence constraints are particularly important for consolidating cross-token visual evidence.

\paragraph{Role of threshold filtering ($p_{\min}$).}
The minimum-mass threshold filters out source chunks whose aggregated attribution mass falls below a predefined proportion,
preventing low-impact assignments from entering the final selection.
This confirms that weak token-level signals—common in visual grounding—
must be explicitly filtered to avoid over-selection of irrelevant image spans.
Threshold filtering therefore acts as a crucial denoising step in the multimodal setting.

